\begin{problem}[33.2]
  determine whether $\{(1, 1, 0), (1, 0, 1), (0, 1, 1)\}$ is a basis for $\R^3$ over $\R$.
\end{problem}

\begin{solution}
  The vectors $(1, 1, 0)$, $(1, 0, 1)$, and $(0, 1, 1)$ are linearly independent. To see this, we can set up the equation:
  \[%
    a(1, 1, 0) + b(1, 0, 1) + c(0, 1, 1) = (0, 0, 0)
  .\]%
  Writing this in matrix form, we have:
  \[%
    \begin{bmatrix}
      1 & 1 & 0 \\
      1 & 0 & 1 \\
      0 & 1 & 1 \\
    \end{bmatrix}
    \begin{bmatrix}
      a \\
      b \\
      c \\
    \end{bmatrix}
    =
    \begin{bmatrix}
      0 \\
      0 \\
      0 \\
    \end{bmatrix}
   .\]%
  Solving this system of equations, we find that the only solution is $a = b = c = 0$. Therefore, the vectors are linearly independent.

  Lastly, we show that the set of vectors spans $\R^3$. To do this, we reduce the matrix to row echelon form:
  \[%
    \begin{bmatrix}
      1 & 1 & 0 \\
      1 & 0 & 1 \\
      0 & 1 & 1 \\
    \end{bmatrix} \leftarrow
    \begin{bmatrix}
      1 & 0 & 0 \\
      1 & 1 & 1 \\
      0 & 1 & 1 \\
    \end{bmatrix}
  .\]%
  Since there's a pivot in every row, the vectors span $\R^3$. Therefore, the set $\{(1, 1, 0), (1, 0, 1), (0, 1, 1)\}$ is a basis for $\R^3$ over $\R$.
\end{solution}

\begin{problem}[33.5]
  Determine if $\{a + b\sqrt{2}~|~a, b \in \Q\}$ is a finite dimensional vector space over $\Q$. If it is, find a basis.
\end{problem}

\begin{solution}
  The set $\{a + b\sqrt{2}~|~a, b \in \Q\}$ is a finite dimensional vector space over $\Q$. To see this, we can express any element of the set as a linear combination of the vectors $1$ and $\sqrt{2}$:
  \[%
    a + b\sqrt{2} = a \cdot 1 + b \cdot \sqrt{2}
  .\]%
  Therefore, the set is spanned by the vectors $1$ and $\sqrt{2}$.

  Next, we check if the vectors $1$ and $\sqrt{2}$ are linearly independent. Suppose there exist rational numbers $x$ and $y$ such that:
  \[%
    x \cdot 1 + y \cdot \sqrt{2} = 0
  .\]%
  This implies that $x + y\sqrt{2} = 0$. Since $\sqrt{2}$ is irrational, the only way for this equation to hold is if $x = 0$ and $y = 0$. Therefore, the vectors are linearly independent.

  Since the vectors $1$ and $\sqrt{2}$ are linearly independent and span the set, they form a basis for the vector space. Thus, a basis for $\{a + b\sqrt{2}~|~a, b \in \Q\}$ over $\Q$ is $\{1, \sqrt{2}\}$.
\end{solution}

\begin{problem}[33.6]
  Determine if $\left.\left\{\frac{a + b\sqrt{3}}{c + d\sqrt{3}}~\right\rvert~a, b, c, d \in \Q \and c + d \ne \sqrt{3}\right\}$ is finite dimensional vector space over $\Q$. If it is, find a basis.
\end{problem}

\begin{solution}
  The set $\left.\left\{\frac{a + b\sqrt{3}}{c + d\sqrt{3}}~\right\rvert~a, b, c, d \in \Q \and c + d \ne \sqrt{3}\right\}$ is a finite dimensional vector space over $\Q$. To see this, we can express any element of the set as a linear combination of the vectors $1$ and $\sqrt{3}$:
  \[%
    \frac{a + b\sqrt{3}}{c + d\sqrt{3}} = \frac{(a + b\sqrt{3})(c - d\sqrt{3})}{(c + d\sqrt{3})(c - d\sqrt{3})} = \frac{(ac - 3bd) + (bc - ad)\sqrt{3}}{c^2 - 3d^2}
  .\]%
  Therefore, the set is spanned by the vectors $1$ and $\sqrt{3}$.

  Next, we check if the vectors $1$ and $\sqrt{3}$ are linearly independent. Suppose there exist rational numbers $x$ and $y$ such that:
  \[%
    x \cdot 1 + y \cdot \sqrt{3} = 0
  .\]%
  This implies that $x + y\sqrt{3} = 0$. Since $\sqrt{3}$ is irrational, the only way for this equation to hold is if $x = 0$ and $y = 0$. Therefore, the vectors are linearly independent.

  Since the vectors $1$ and $\sqrt{3}$ are linearly independent and span the set, they form a basis for the vector space. Thus, a basis for $\left.\left\{\frac{a + b\sqrt{3}}{c + d\sqrt{3}}~\right\rvert~a, b, c, d \in \Q \and c + d \ne \sqrt{3}\right\}$ over $\Q$ is $\{1, \sqrt{3}\}$.
\end{solution}

\begin{problem}[33.7]
  Determine if $\R$ is a finite dimensional vector space over $\{a + b\sqrt{2}~|~a, b \in \Q\}$. If it is, find a basis.
\end{problem}

\begin{solution}
  The set $\R$ is not a finite dimensional vector space over $\{a + b\sqrt{2}~|~a, b \in \Q\}$. To see this, we can consider the dimension of $\R$ over $\Q$. Since $\R$ is an infinite-dimensional vector space over $\Q$, it cannot be a finite-dimensional vector space over any subfield of $\R$ that contains $\Q$. Therefore, $\R$ is not a finite dimensional vector space over $\{a + b\sqrt{2}~|~a, b \in \Q\}$.
\end{solution}

\begin{problem}[33.8]
  Determine if $\R$ is a finite dimensional vector space over $\C$. If it is, find a basis.
\end{problem}

\begin{solution}
  The set $\R$ is not a vector space over $\C$. For a set $V$ to be a vector space over a field $F$, it must be closed under scalar multiplication, meaning $a \cdot v \in V$ for all $a \in F$ and $v \in V$. If we take the scalar $i \in \C$ and the vector $1 \in \R$, their product is $i \cdot 1 = i$. Since $i \notin \R$, the closure property fails. Therefore, $\R$ cannot be a vector space over $\C$.
\end{solution}

\begin{problem}[33.9]
  Determine if $\R$ is a finite dimensional vector space over $\Q$. If it is, find a basis.
\end{problem}

\begin{solution}
  The set $\R$ is not a finite dimensional vector space over $\Q$. To see this, we can consider the dimension of $\R$ over $\Q$. Since $\R$ is an infinite-dimensional vector space over $\Q$, it cannot be a finite-dimensional vector space over $\Q$. Therefore, $\R$ is not a finite dimensional vector space over $\Q$.
\end{solution}

\begin{problem}[33.10]
  There is a field $E$ with 32 elements. Determine which prime field is isomorphic with a subfield of $E$ and determine the dimension of $E$ over its prime field.
\end{problem}

\begin{solution}
  To determine which prime field is isomorphic with a subfield of $E$, we need to find the characteristic of $E$. The characteristic of a field is the smallest positive integer $p$ such that $p \cdot 1 = 0$ in the field. Since $E$ has 32 elements, it must have characteristic 2, because 32 is a power of 2. Therefore, the prime field that is isomorphic with a subfield of $E$ is $\F_2$, the field with 2 elements.
\end{solution}

\begin{problem}[33.21]
  Prove that if $V$ is a finite-dimensional vector space over a field $F$, then a subset $\{\beta_1,$ $\beta_2,$ $\cdots,$ $\beta_n\}$ of $V$ is a basis for $V$ over $F$ if and only if every vector in $V$ can be expressed uniquely as a linear combination of the $\beta_i$.
\end{problem}

\begin{solution}
  Assume that $\{\beta_1, \beta_2, \cdots, \beta_n\}$ is a basis for $V$ over $F$. By definition of a basis, the vectors are linearly independent and span $V$. Therefore, every vector in $V$ can be expressed as a linear combination of the $\beta_i$. To show uniqueness, suppose there are two different linear combinations of the $\beta_i$ that represent the same vector:
  \[%
    a_1\beta_1 + a_2\beta_2 + \cdots + a_n\beta_n = b_1\beta_1 + b_2\beta_2 + \cdots + b_n\beta_n
  .\]%
  Rearranging this equation gives:
  \[%
    (a_1 - b_1)\beta_1 + (a_2 - b_2)\beta_2 + \cdots + (a_n - b_n)\beta_n = 0
  .\]%
  Since the $\beta_i$ are linearly independent, it follows that $a_i - b_i = 0$ for all $i$, which means $a_i = b_i$ for all $i$. Thus, the representation of each vector in terms of the basis is unique.

  Now, conversely, assume assume that every vector in $V$ can be expressed uniquely as a linear combination of $\{\beta_1, \beta_2, \cdots, \beta_n\}$. To show that these vectors form a basis, we need to show that they are linearly independent and span $V$.

  First, we show linear independence. Suppose there exist scalars $c_1, c_2, \cdots, c_n$ such that:
  \[%
    c_1\beta_1 + c_2\beta_2 + \cdots + c_n\beta_n = 0
  .\]%
  Since the representation of each vector is unique, the only way for this equation to hold is if $c_i = 0$ for all $i$. Therefore, the vectors are linearly independent.

  Next, we show that the vectors span $V$. Since every vector in $V$ can be expressed as a linear combination of $\{\beta_1, \beta_2, \cdots, \beta_n\}$, it follows that the vectors span $V$. Therefore, $\{\beta_1, \beta_2, \cdots, \beta_n\}$ is a basis for $V$ over $F$.
\end{solution}

\begin{problem}[33.23]
  Prove that every finite-dimensional vector space $V$ of dimension $n$ over a field $F$ is isomorphic to the vector space $F^n$ of Exercise 19.
\end{problem}

\begin{solution}
  Let $V$ be a finite-dimensional vector space of dimension $n$ over a field $F$. Since $V$ has dimension $n$, there exists a basis $\{\beta_1, \beta_2, \cdots, \beta_n\}$ for $V$ over $F$. We can define a function $\phi : V \to F^n$ as follows: For any vector $\alpha \in V$, express $\alpha$ as a linear combination of the basis vectors:
  \[%
    \alpha = a_1\beta_1 + a_2\beta_2 + \cdots + a_n\beta_n
  .\]%
  Then we define:
  \[%
    \phi(\alpha) = (a_1, a_2, \cdots, a_n)
  .\]%
  This function $\phi$ is well-defined because the representation of each vector in terms of the basis is unique.

  Next, we need to show that $\phi$ is an isomorphism. To show that $\phi$ is linear, we check the following properties for any vectors $\alpha, \gamma \in V$ and any scalar $a \in F$:
  \begin{align*}
    \phi(\alpha + \gamma) &= \phi((a_1\beta_1 + a_2\beta_2 + \cdots + a_n\beta_n) + (c_1\beta_1 + c_2\beta_2 + \cdots + c_n\beta_n)) \\
                          &= (a_1 + c_1, a_2 + c_2, \cdots, a_n + c_n) = (a_1, a_2, \cdots, a_n) + (c_1, c_2, \cdots, c_n) = \phi(\alpha) + \phi(\gamma)
  ,\end{align*}
  and
  \begin{align*}
    \phi(a\alpha) &= \phi(a(a_1\beta_1 + a_2\beta_2 + \cdots + a_n\beta_n)) = (aa_1, aa_2, \cdots, aa_n) = a(a_1, a_2, \cdots, a_n) = a\phi(\alpha)
  .\end{align*}
  Therefore, $\phi$ is a linear transformation.

  To show that $\phi$ is injective, suppose $\phi(\alpha) = \phi(\gamma)$ for some $\alpha, \gamma \in V$. This means that the coordinates of $\alpha$ and $\gamma$ with respect to the basis are the same, which implies that $\alpha = \gamma$. Therefore, $\phi$ is injective.

  To show that $\phi$ is surjective, let $(b_1, b_2, \cdots, b_n) \in F^n$ be any vector. We can construct a vector $\alpha \in V$ as follows:
  \[%
    \alpha = b_1\beta_1 + b_2\beta_2 + \cdots + b_n\beta_n
  .\]%
  Then we have:
  \[%
    \phi(\alpha) = (b_1, b_2, \cdots, b_n)
  .\]%
  Therefore, $\phi$ is surjective.

  Since $\phi$ is a linear transformation that is both injective and surjective, it is an isomorphism. Thus, $V$ is isomorphic to $F^n$.
\end{solution}

\begin{problem}[33.24]
  Let $V$ and $V'$ be vector spaces over the same field $F$. A function $\phi : V \to V'$ is a \textit{linear transformation} of $V$ into $V'$ if the following conditions are satisfied for all $\alpha, \beta \in V$, and $a \in F$:
  \begin{gather*}
    \phi(\alpha + \beta) = \phi(\alpha) + \phi(\beta) \\
    \phi(a\alpha) = a\phi(\alpha)
  \end{gather*}
  \begin{enumerate}
    \item If $\{\beta_i | i \in I\}$ is a basis for $V$ over $F$, show that a linear transformation $\phi : V \to V'$ is completely determined by the vectors $\phi(\beta_i) \in V'$.
    \item Let $\{\beta_i | i \in I\}$ be a basis for $V$, and let $\{\beta_i' | i \in I\}$ be any set of vectors, not necessarily distinct, of $V'$. Show that there exists exactly one linear transformation $\phi : V \to V'$ such that $\phi(\beta_i) = \beta_i'$.
  \end{enumerate}
\end{problem}

\begin{solution}[(i)]
  Assume that $\{\beta_i | i \in I\}$ is a basis for $V$ over $F$. Since $\{\beta_i | i \in I\}$ is a basis, every vector $\alpha \in V$ can be expressed as a linear combination of the basis vectors:
  \[%
    \alpha = a_1\beta_{i_1} + a_2\beta_{i_2} + \cdots + a_k\beta_{i_k}
  .\]%
  Applying the linear transformation $\phi$ to both sides, we have:
  \[%
    \phi(\alpha) = \phi(a_1\beta_{i_1} + a_2\beta_{i_2} + \cdots + a_k\beta_{i_k}) = a_1\phi(\beta_{i_1}) + a_2\phi(\beta_{i_2}) + \cdots + a_k\phi(\beta_{i_k})
  .\]%
  Therefore, the value of $\phi(\alpha)$ is completely determined by the values of $\phi(\beta_i)$ for the basis vectors $\beta_i$. Since every vector in $V$ can be expressed as a linear combination of the basis vectors, and the linear transformation is determined by its action on the basis vectors, it follows that $\phi$ is completely determined by the vectors $\phi(\beta_i) \in V'$.
\end{solution}

\begin{solution}[(ii)]
  Let $\{\beta_i | i \in I\}$ be a basis for $V$, and let $\{\beta_i' | i \in I\}$ be any set of vectors in $V'$. We want to show that there exists exactly one linear transformation $\phi : V \to V'$ such that $\phi(\beta_i) = \beta_i'$ for all $i \in I$.

  First, we define a function $\phi : V \to V'$ as follows: For any vector $\alpha \in V$, express $\alpha$ as a linear combination of the basis vectors:
  \[%
    \alpha = a_1\beta_{i_1} + a_2\beta_{i_2} + \cdots + a_k\beta_{i_k}
  .\]%
  Then we define:
  \[%
    \phi(\alpha) = a_1\beta_{i_1}' + a_2\beta_{i_2}' + \cdots + a_k\beta_{i_k}'
  .\]%
  This definition ensures that $\phi(\beta_i) = \beta_i'$ for all $i \in I$.

  Next, we need to show that $\phi$ is a linear transformation. For any vectors $\alpha, \gamma \in V$ and any scalar $a \in F$, we have:
  \begin{align*}
    \phi(\alpha + \gamma) &= \phi((a_1\beta_{i_1} + a_2\beta_{i_2} + \cdots + a_k\beta_{i_k}) + (c_1\beta_{j_1} + c_2\beta_{j_2} + \cdots + c_m\beta_{j_m})) \\
                          &= (a_1 + c_1)\beta_{i_1}' + (a_2 + c_2)\beta_{i_2}' + \cdots + (a_k + c_k)\beta_{i_k}' + (c_1)\beta_{j_1}' + (c_2)\beta_{j_2}' + \cdots + (c_m)\beta_{j_m}' \\
                          &= (a_1\beta_{i_1}' + a_2\beta_{i_2}' + \cdots + a_k\beta_{i_k}') + (c_1\beta_{j_1}' + c_2\beta_{j_2}' + \cdots + c_m\beta_{j_m}') \\
                          &= \phi(\alpha) + \phi(\gamma)
  ,\end{align*}
  and
  \begin{align*}
    \phi(a\alpha) &= \phi(a(a_1\beta_{i_1} + a_2\beta_{i_2} + \cdots + a_k\beta_{i_k})) \\
                  &= (aa_1)\beta_{i_1}' + (aa_2)\beta_{i_2}' + \cdots + (aa_k)\beta_{i_k}' \\
                  &= a(a_1\beta_{i_1}' + a_2\beta_{i_2}' + \cdots + a_k\beta_{i_k}') \\
                  &= a\phi(\alpha)
  .\end{align*}
  Therefore, $\phi$ is a linear transformation.
\end{solution}
